\section{GraphDrone.fit() Phase I-B Expert Factory}

Phase I-B upgrades GraphDrone from a manifest-only bootstrap package to a package that can construct a portfolio of experts directly from training data. The core shift is that \texttt{GraphDrone.fit()} now accepts an explicit set of expert build specifications:
\[
\mathcal{E} = \{(d_v, a_v, m_v, \theta_v)\}_{v=1}^V,
\]
where \(d_v\) is the typed descriptor, \(a_v\) is the input adapter, \(m_v\) is the expert model family, and \(\theta_v\) are model parameters.

For each expert \(v\), Phase I-B fits
\[
f_v \leftarrow \mathrm{FitExpert}(a_v(X_{\mathrm{train}}), y_{\mathrm{train}}; m_v, \theta_v),
\]
then packages the fitted experts into a validated portfolio with a required anchor \(\mathrm{FULL}\). This means benchmark scripts no longer need to fit TabPFN experts inline; they only define an expert plan and call the package API.

The benchmark adapter remains intentionally thin. It constructs a \texttt{GraphDroneBenchmarkExpertPlan} from registered OpenML split metadata and view declarations, then hands those specs to the model package. In the current phase, supported expert kinds are \texttt{constant}, \texttt{linear}, and \texttt{tabpfn\_regressor}. The important architectural result is not model novelty but ownership transfer: expert construction now lives inside \texttt{src/graphdrone\_fit/}.

Validation for this phase consisted of:
\begin{itemize}
\item targeted unit and adapter tests, passing at 12 tests,
\item a real smoke benchmark on \texttt{houses},
\item external Gemini review of the file-level implementation.
\end{itemize}

The smoke run behaved as expected for a bootstrap-only router: the \texttt{FULL} expert and the explicit \texttt{GraphDroneFit\_bootstrap} output matched, while specialist experts were still exposed separately for inspection. This confirms that the package can now own expert construction without yet claiming contextual routing gains.

The remaining gaps are operational rather than conceptual: serialization of spec-built portfolios, multi-device validation of the new expert-construction path, and replacement of the bootstrap router with a real contextual set-router in later phases.
